\section{Background}
\label{sec:background}

\subsection{Urban morphometrics}
\label{sec:bg-morphometrics}
Over the past decade the measurement of urban form has consolidated into a
reproducible, quantitative science. Quantitative morphometric descriptors lay the
mathematical foundations for analyzing the evolution of urban patterns and their
correlation with neighborhood vibrancy \citep{dibble2019OriginSpaces,
fleischmann2021MeasurinUrban, lu2019ExplorinAssociat}. Morphological tessellation
provides a consistent, cadastre-independent proxy for the plot, the elementary
unit at which built form is generated \citep{fleischmann2020Morphologic}, and a
numerical taxonomy built on contextual characters and unsupervised classification
turns the description of fabric into a repeatable procedure
\citep{fleischmann2022MethodolFoundati}. The approach now scales: street-based
analyses support nationwide fabric taxonomies \citep{araldi2025MultiLevel}, and
unsupervised clustering of satellite-derived morphology yields a global typology
of the urban fabric \citep{debray2025UniversalPatterns}. Recent applications have
mapped these typomorphological profiles to COVID-19 dynamics, Airbnb distribution,
and socioeconomic deprivation \citep{venerandi2023UrbanForm, venerandi2023UrbanFormb,
venerandi2024UrbanForm}, while neighborhood delineation techniques combine morphology
with flow datasets \citep{govind2024DelineatNeighbor}. LiDAR-based multi-scale
aerodynamic roughness modeling and historical heritage-area vitality assessments
further enrich this morphometric toolkit \citep{an2024StudyMorphome,
wu2022CultivatHistoric}. Interpretable multi-indicator pipelines
\citep{vartholomaios2025Detection} and deep-learning classifiers working from
imagery and figure-ground maps \citep{fang2024Urbanclassifier,
wangJ2024FigureGround} extend these methodologies. The decisive feature of this
literature is also its limitation: it is predominantly \emph{descriptive} and
form-only: it names fabric types with rigour but rarely connects them to the
outcomes those types produce.

\subsection{Climate hazards}
\label{sec:bg-hazard}
A second, largely separate literature establishes that urban morphology is a
control on climate hazard rather than a correlate of it. Form shapes land-surface
temperature through marginal and interaction effects that vary by season and time
of day \citep{wangZ2025MarginalInteraction}, with density and Spacematrix-style
morphometrics exercising a key control on local microclimates \citep{li2020InfluencDensity}.
Heat and pluvial flooding increasingly behave as compound, interacting
hazards \citep{imroz2025IntegratedAssess}, requiring integrated flood vulnerability
index formulations for coastal cities \citep{balica2012FloodVulnerab}. Quantitative
studies examine resilient flood risk management strategies in deltaic settings
\citep{chan2018TowardsResilien}, the role of nature-based solutions in building
urban resilience \citep{bush2019BuildingUrban}, and how local flooding events can
induce abrupt, cascading failures of street networks \citep{wang2019LocalFloods}.
Understanding the homogeneous and heterogeneous spatial patterns of diurnal and
nocturnal heat hotspots is also essential for synergetic mitigation planning
\citep{liu2025HomogeneHeteroge}. The cooling delivered by green and blue
nature-based solutions is now quantified across morphological types and spatial
scales \citep{wei2025NatureBased}. Methodologically, explainable machine learning
Gradient boosting with SHAP attribution has become the standard lens on the
morphology$\rightarrow$temperature relationship, exposing non-linear and
spatially heterogeneous effects that linear models miss
\citep{wangZ2025PlainPlateau, liu2026Nonlinear, li2026BlockGrid}. We adopt this
lens but relocate it from the grid or Local Climate Zone tile
\citep{stewart2012LocalClimate} to the street/cell scale, and carry its mechanism attributions into the interpretation of a typology and a
decision step, rather than stopping at explanation. The gap this literature leaves is that it is
typically single-hazard, resolved at coarse units, and decoupled from
accessibility, equity and any downstream prioritization.

\subsection{Street networks}
\label{sec:bg-network}
The street network is both a morphological object and the substrate of everyday
access. Large-scale comparative street network modeling and studies of network
planarity show how structural configurations differ across global regions
\citep{boeing2020MultiScale, boeing2020PlanaritStreet}. Reproducible open tooling
has made network retrieval, modelling and indicator computation routine and
comparable across cities \citep{boeing2017Osmnx, boeing2022StreetNetwork,
boeing2025ModelingNetworks}, while street-based pedestrian-perspective models
bridge local form with regional scales \citep{araldi2019FromStreet}. Configurational
reasoning in the space-syntax tradition links network structure to movement and
co-presence \citep{hillier1984SocialLogic}, albeit with known limits at fine
urban-design scales \citep{pafka2020LimitsSpace}. The proximity turn, including
assessment frameworks for the 15-minute city, ties urban form directly to accessibility
outcomes \citep{wangH2025FifteenMinute}. Here too, however, accessibility is
usually treated apart from hazard: walkable does not mean cool, and a
well-connected fabric can still be a heat or flood trap. Our workflow keeps
movement potential and hazard exposure in the same unit so that the two can be
read against each other.

\subsection{Climate justice}
\label{sec:bg-vulnerability}
Exposure is not risk. The social-vulnerability tradition, from the foundational
social-vulnerability index \citep{cutter2003SocialVulnerab} to flood-specific
constructions \citep{aroca2017ConstrucIntegrat} and heat-coupled indices
\citep{sabrin2020DevelopiVulnerab}, insists that the capacity to cope is unevenly
distributed. Recent work makes this operational and transferable at scale. An
open, validated global urban heat-vulnerability index
\citep{turner2025Guhvi}, while empirical results show that perceived heat stress and
demographic vulnerability are frequently decoupled across urban forms
\citep{iqbal2025HeatStress}. Integrating vulnerability into adaptation is
increasingly framed not as an add-on but as a justice imperative
\citep{neumann2026ClimateJustice}. No existing framework, however, binds this vulnerability axis to the
\emph{form} that produces differential exposure in the first place.

\subsection{Prioritization}
\label{sec:bg-mcda}
Finally, adaptation assessment must convert indicators into priorities.
Multi-criteria decision analysis offers transparent, auditable aggregation. The
technique for order of preference by similarity to ideal solution (TOPSIS) is a
well-established example \citep{hwang1981MultipleAttribute}, while entropy
weighting reduces analyst discretion. But composite indices collapse genuine
trade-offs into a single score. Multi-objective (Pareto) optimization, already
established for urban-heat and green-infrastructure trade-offs
\citep{zhang2024MultiObjective, zhu2025SpongeCities} and for spatial conservation
prioritization \citep{schroter2016SpatialPrioriti}, instead separates
unambiguously dominated options (clear intervention candidates) from non-dominated
trade-off options. It has not, however, been applied to fabric-level adaptation
priority. We pair a Pareto frontier with a TOPSIS ranking and Monte-Carlo weight
perturbation, so that the priority output is both decision-ready and explicit
about its own robustness.

These five strands are individually mature but rarely meet. No existing workflow
binds reproducible morphometrics, multi-hazard exposure and social vulnerability
on a sampled metropolitan grid within one commensurable unit, and then converts
them into explainable, Pareto-aware adaptation priorities. This study occupies
that gap with a single auditable unit that is at once morphometric,
configurational, hazard-aware and equity-aware within an explain-then-optimize
pipeline.
